\documentclass[11pt,a4paper]{article}
\usepackage[margin=25mm]{geometry}
\usepackage[T1]{fontenc}
\usepackage{lmodern}
\usepackage{amsmath,amssymb}
\usepackage{booktabs,longtable,array}
\usepackage{graphicx}
\usepackage{hyperref}
\hypersetup{colorlinks=true,linkcolor=blue,urlcolor=blue}
\setlength{\parskip}{0.5em}
\setlength{\parindent}{0pt}
\setlength{\emergencystretch}{3em}
\newcommand{\starred}[1]{#1^{*}}
\newcommand{\NAIRU}{u^{*}}
\newcommand{\gap}{g}
\newcommand{\eps}{\varepsilon}

\title{Korean SMOG Model\\Coefficient and Fit Analysis}
\author{Model output review}
\date{16 September 2026\\Saved model results through 2026Q1}

\begin{document}
\maketitle

\section{Scope and interpretation}
This report analyzes the standalone Korean Python SMOG model.
Estimates come from its saved
parameter JSON and accepted-iteration likelihood history. The generation
sample is 2000Q1--2022Q4 (92 quarters). The update applies those fixed
coefficients to later data; the latest complete saved result is 2026Q1.
The latent potential output and NAIRU are model estimates, not directly
observed quantities.

\section{Results and charts}
The fixed-parameter update runs through 2026Q1. Its saved smoothed output
gap is $-0.052\%$ of potential output and its NAIRU is $2.761\%$.
Observed unemployment is $2.867\%$, or 0.106 percentage points above the
fitted NAIRU. Smoothed states use the full updated sample; filtered states
use observations available through each plotted quarter, with the same
historical coefficients. The charts are model estimates, not direct
measurements of potential output or the NAIRU.

\begin{figure}[p]
\centering
\includegraphics[width=0.85\linewidth,height=0.28\textheight,keepaspectratio]{\detokenize{../results/figures/kr_smog_output_gap_smoothed.png}}
\caption{Smoothed output gap.}
\end{figure}

\begin{figure}[p]
\centering
\includegraphics[width=0.85\linewidth,height=0.28\textheight,keepaspectratio]{\detokenize{../results/figures/kr_smog_output_gap_filtered.png}}
\caption{Filtered output gap.}
\end{figure}

\begin{figure}[p]
\centering
\includegraphics[width=0.85\linewidth,height=0.28\textheight,keepaspectratio]{\detokenize{../results/figures/kr_smog_unemployment_and_nairu_smoothed.png}}
\caption{Observed unemployment and smoothed NAIRU.}
\end{figure}

\begin{figure}[p]
\centering
\includegraphics[width=0.85\linewidth,height=0.28\textheight,keepaspectratio]{\detokenize{../results/figures/kr_smog_unemployment_and_nairu_filtered.png}}
\caption{Observed unemployment and filtered NAIRU.}
\end{figure}

\begin{figure}[p]
\centering
\includegraphics[width=0.85\linewidth,height=0.28\textheight,keepaspectratio]{\detokenize{../results/figures/kr_smog_unemployment_gap_smoothed.png}}
\caption{Smoothed NAIRU minus observed unemployment.}
\end{figure}

\begin{figure}[p]
\centering
\includegraphics[width=0.85\linewidth,height=0.28\textheight,keepaspectratio]{\detokenize{../results/figures/kr_smog_unemployment_gap_filtered.png}}
\caption{Filtered NAIRU minus observed unemployment.}
\end{figure}

\clearpage

\section{Reading the current output gap}
The latest saved Korean gap is $-0.052\%$ in 2026Q1, quite close to zero;
observed unemployment is 0.106 percentage points above the fitted NAIRU.
The local core-inflation reading is $2.176\%$ and the expectation input is
$2.633\%$. The smoothed gap contributes just $-0.006$ percentage points to
the fitted inflation equation. Its full conditional fitted value is
$2.153\%$, leaving an actual-minus-fitted difference of $0.022$ points.
This uses the smoothed state and is not an out-of-sample forecast.

The evidence here is weak for calling the output gap materially too negative:
its estimated level and direct inflation contribution are both small. The
unemployment equation helps locate the latent gap, but lagged inflation,
expectations, import prices, the wage signal, and measurement noise also
shape the joint fit. Current inflation alone cannot identify which latent
state or coefficient should change.

\section{State-space structure}
The model uses seven states,
$s_t=(\gap_t,\gap_{t-1},\gap_{t-2},y^*_t,y^*_{t-1},\NAIRU_t,\NAIRU_{t-1})'$, where
$y_t$ is log output, $\gap_t=y_t-y_t^*$ is a log output gap, and $\NAIRU_t$
is the latent structural unemployment rate. The reported output gap is
$100\gap_t$, in percent of potential output. The transition equations
are
\begin{align}
 \gap_t &= \phi_1\gap_{t-1}+\phi_2\gap_{t-2}+\eta^g_t,\\
 y_t^* &= y_{t-1}^*+g^*+\eta^y_t,\\
 \NAIRU_t &= \NAIRU_{t-1}+\eta^u_t.
\end{align}
The output gap follows an AR(2), potential log output is a random walk with
quarterly drift $g^*$, and the NAIRU is a random walk. The three state shocks
are independent in the implemented diagonal covariance matrix. The state
vector also carries lags needed by the observation equations.

At the beginning of each generation run, HP-based gap, potential output, and
NAIRU values set a known initial state with zero prior covariance. The model reads its historical HP columns from the Generation tab. The HP series initialize the
filter and provide regression-based optimizer starting values; the saved
maximum-likelihood coefficients are fitted by Python from the country's own
signals. This fixed initialization can matter for early likelihood terms and
the resulting parameter estimates.

\section{Korea: equations and estimates}
Korea has four observed signals: log GDP $y_t$, unemployment $u_t$, core
inflation $\pi_t$, and nominal unit labour cost growth $w_t$. Here
$\pi_t^e$ is the Bank of Korea expectation input and $p_{t-1}^m$ is the
lagged four-quarter import-price change. The implemented observation
equations are
\begin{align}
 y_t &= y_t^*+\gap_t+\eps^y_t,\\
 u_t &= \NAIRU_t+\lambda_3\gap_t
       +\rho_1(u_{t-1}-\NAIRU_{t-1})+\eps^u_t,\\
 \pi_t &= (1-\beta_1-\beta_2)\pi_t^e
      +\beta_1\pi_{t-1}+\beta_2\pi_{t-2}
      +\lambda_1\gap_t+\psi p_{t-1}^m+\eps^\pi_t,\\
 w_t &= (1-\beta_4)\pi_t^e+\beta_4\pi_{t-1}
      +\lambda_2(u_{t-1}-\NAIRU_{t-1})+\eps^w_t.
\end{align}
The wage equation connects labour-market slack to unit labour cost growth.
Because the gap is in log units, a one percentage point change in the output
gap multiplies $\lambda_1$ or $\lambda_3$ by $0.01$.

\begin{longtable}{@{}>{$}l<{$}r p{0.65\linewidth}@{}}
\caption{Korean maximum-likelihood coefficients. Values are rounded.}\\
\toprule
\text{Parameter} & \multicolumn{1}{c}{\text{Estimate}} & Meaning in the implemented model\\
\midrule
\endfirsthead
\toprule
\text{Parameter} & \multicolumn{1}{c}{\text{Estimate}} & Meaning\\
\midrule
\endhead
\beta_1 & 1.311422 & Weight on one-quarter lagged core inflation in the inflation equation.\\
\beta_2 & -0.330134 & Weight on two-quarter lagged core inflation.\\
\beta_4 & 1.966103 & Weight on lagged inflation in the wage equation; the implied expectation weight is $-0.9661$.\\
\lambda_1 & 11.745785 & Inflation loading on the log output gap; a one percentage point gap gives about 0.1175 percentage points through this term.\\
\lambda_2 & -8.279811 & Wage loading on lagged unemployment minus lagged NAIRU; a one percentage point larger unemployment gap lowers predicted wage growth by about 8.28 percentage points, other terms fixed.\\
\lambda_3 & -8.620103 & Unemployment loading on the log output gap; a one percentage point output gap gives about $-0.0862$ percentage points.\\
\rho_1 & 0.889138 & Persistence of lagged unemployment relative to lagged NAIRU.\\
\phi_1 & 1.290464 & First autoregressive coefficient of the latent output gap.\\
\phi_2 & -0.596930 & Second autoregressive coefficient of the latent output gap.\\
\psi & 0.002122 & Loading of lagged four-quarter import-price change in inflation.\\
g^* & 0.009052 & Quarterly drift of potential log output, about 3.62\% at an annualized log-growth rate.\\
\eps_1 & 0.005770 & Standard deviation of log GDP measurement noise; about 0.577\% in log-output terms.\\
\eps_2 & 0.074402 & Standard deviation of unemployment measurement noise, in percentage points.\\
\eps_3 & 0.285165 & Standard deviation of core-inflation measurement noise, in percentage points.\\
\eps_4 & 4.895773 & Standard deviation of unit labour cost growth measurement noise, in percentage points.\\
\eps_5 & 0.043368 & Code-scale parameter for the gap shock: actual shock standard deviation is $\eps_5/10=0.004337$ log units, or about 0.434\%.\\
\eps_6 & 0.006277 & Standard deviation of the potential-log-output state shock.\\
\eps_7 & 0.157126 & Standard deviation of the quarterly NAIRU state shock, in unemployment percentage points.\\
\bottomrule
\end{longtable}

The Korean inflation expectation weight is $1-\beta_1-\beta_2=0.0187$.
The large wage measurement standard deviation, $\eps_4$, means that the
fourth signal is comparatively noisy in its own units. It contributes to the
likelihood, but an individual wage observation need not be close to its
conditional fitted value. The sizeable negative $\lambda_2$ is a partial
effect on the wage equation, conditional on its lagged inflation and
expectation terms; it should not be read as a standalone causal estimate.

\section{Likelihood fit and convergence}
The script uses SciPy's bounded L-BFGS-B optimizer on the Kalman-filter log
likelihood. Starting guesses come from Korean historical data. The
accepted-iteration history stores all 18 parameters, the full-sample log
likelihood, and quarter-level contributions. Trial AR(2) coefficients
outside the implemented stability inequalities are penalized.

\begin{center}
\begin{tabular}{@{}lr@{}}
\toprule
Fit statistic & Korea\\
\midrule
Historical quarters & 92\\
Accepted iterations & 129\\
Initial log likelihood & $-33{,}775.3474$\\
Iteration 25 & 12.1332\\
Iteration 50 & 12.3798\\
Final log likelihood & 12.3840\\
\bottomrule
\end{tabular}
\end{center}

The initial likelihood is extremely poor, but accepted steps increased it
monotonically. The last history row agrees with the saved parameter JSON.
Numerical convergence from these starting values and bounds does not
establish a global optimum, tight identification, or forecast accuracy. The
script does not export standard errors or out-of-sample forecast scores.
Quarter-level likelihood contributions are Kalman one-step-ahead predictive
density terms, not squared-error metrics. Early missing regressors cause the
affected signals to be masked.

\section{Reproduction and source files}
From this model folder, run the generation script to re-estimate historical
coefficients and overwrite the history CSVs. Then run the update script to
apply the fixed fit to later quarters. Source files used for this report:
\begin{itemize}
\item Generation and update scripts in the \texttt{code} folder.
\item Fitted parameter JSON at the model-folder root.
\item The fit-history and likelihood-by-quarter CSVs in \texttt{output}.
\item The updated state CSV in \texttt{results} and the \texttt{Update}
      tab of the local model-input workbook.
\end{itemize}
Numeric tables are a snapshot and do not update automatically when Python is
rerun.

\end{document}
